\chapter{Finite zero-mode supertrace и знак Coleman--Weinberg коэффициента}

Quantum scale gate требует вычислить
\begin{equation}
 B=\frac1{64\pi^2\chi^4}\Str M^4(\chi)
\end{equation}
для заранее объявленного поля. Первый контролируемый слой --- конечные
zero-modes без KK-башен и gauge/ghost completion.

\section{Trace-кинетическая метрика}

На quotient по непPhysical rephasing используем координаты
\begin{equation}
 x=re^{i\phi/2},
 \qquad
 z=se^{i\phi/2}.
\end{equation}
Канонический Hilbert--Schmidt kinetic form задаётся
\begin{equation}
 ds^2_{\mathrm{kin}}
 =\Tr_{\mathrm{orb}}(dD_F\,dD_F).
\end{equation}
Прямое вычисление даёт
\begin{equation}
 G(r,s,\phi)
 =
 \begin{pmatrix}
 4&0&0\\
 0&4&0\\
 0&0&r^2+s^2
 \end{pmatrix}.
\end{equation}
Во flattening-vacuum
\begin{equation}
 r=s=\frac1{\sqrt2},
 \qquad
 \phi=\frac{\pi}{2}
\end{equation}
получаем
\begin{equation}
 G_*=\operatorname{diag}(4,4,1).
\end{equation}

\section{Канонические scalar masses}

Hessian flattening-потенциала в тех же координатах равен
\begin{equation}
 H_*=\operatorname{diag}(32,32,8).
\end{equation}
Физические массы определяются обобщённой задачей
\begin{equation}
 H_*v=m^2G_*v.
\end{equation}
Поэтому три quotient-скаляра точно вырождены:
\begin{equation}
 m_{s,1}^2=m_{s,2}^2=m_{s,3}^2=8\chi^2.
\end{equation}

\begin{proposition}[Canonical finite-scalar spectrum]
При общей trace-нормировке kinetic и flattening terms конечный
Dirac-модуль содержит три real scalar zero-modes с одинаковой
безразмерной массой \(m_s^2/\chi^2=8\).
\end{proposition}

\section{Fermion ledger}

В минимуме
\begin{equation}
 B B^\dagger=\chi^2\mathbf1_2,
\end{equation}
поэтому chiral block имеет два единичных сингулярных значения. Минимальное
четырёхмерное чтение соответствует двум Dirac-парам массы
\begin{equation}
 m_{f,1}=m_{f,2}=\chi.
\end{equation}
KO6 particle--antiparticle doubling не удваивает этот physical count, если
используется ранее сформулированный orbit half-trace.

\section{Минимальный supertrace}

Для real scalar вес равен \(+1\), для Dirac fermion --- \(-4\). Поэтому
\begin{align}
 \frac{\Str M^4}{\chi^4}
 &=3\cdot8^2-2\cdot4\cdot1^4\\
 &=192-8\\
 &=184=8\cdot23.
\end{align}
Минимальный finite zero-mode coefficient равен
\begin{equation}
 \boxed{
 B_0^{\mathrm{fin}}
 =\frac{184}{64\pi^2}
 =\frac{23}{8\pi^2}>0.}
\end{equation}

\begin{theorem}[Conditional finite-sector \(B>0\)]
При канонической общей trace-нормировке, трёх propagating scalar modes и
двух physical Dirac-парах конечный zero-mode сектор даёт положительный
Coleman--Weinberg коэффициент \(23/(8\pi^2)\).
\end{theorem}

Число \(23\) здесь является арифметическим результатом supertrace ledger,
а не автоматически тем же rank-complement, что в нейтринном секторе.
Идентификация этих двух появлений запрещена без общей representation
теоремы.

\section{Relative kinetic-normalization red-team}

Пусть kinetic term имеет относительный коэффициент \(\kappa>0\):
\begin{equation}
 \mathcal L_{\mathrm{kin}}
 =\frac{\kappa}{2}
 \Tr_{\mathrm{orb}}(\partial_\mu D_F\partial^\mu D_F).
\end{equation}
После канонической нормировки
\begin{equation}
 m_s^2=\frac8\kappa\chi^2,
\end{equation}
и supertrace становится
\begin{equation}
 \frac{\Str M^4}{\chi^4}
 =\frac{192}{\kappa^2}-8.
\end{equation}
Знак положителен только при
\begin{equation}
 \kappa<\sqrt{24}.
\end{equation}
Следовательно, \(\kappa=1\) не может считаться безусловным следствием
одного лишь finite-module trace: необходимо вывести относительный
коэффициент spacetime kinetic term и algebraic flattening potential.

\begin{nogocriterion}[Finite \(B\)-normalization gap]
Положительное значение \(B_0^{\mathrm{fin}}=23/(8\pi^2)\) повышается до
теоремы parent action только если одна spectral/variational construction
фиксирует \(\kappa=1\), fermion Yukawa normalization и orbit half-trace.
Иначе знак \(B\) зависит от непрерывной kinetic/potential normalization.
\end{nogocriterion}

\section{Неполнота полного spectrum ledger}

В текущий \(B_0^{\mathrm{fin}}\) не включены:
\begin{enumerate}
 \item возможная gauge-мода конечной алгебры и её ghost complex;
 \item dilaton/radion fluctuation;
 \item nonzero KK levels на \(K\);
 \item scheme-dependent subtraction полного tower supertrace.
\end{enumerate}
Поэтому найден положительный finite seed, но не полный \(B\).

\section{Следующий гейт}

Следует вывести \(\kappa\) из product spectral action и затем вычислить
regulated разность
\begin{equation}
 \Delta B_{\mathrm{KK+gauge}}
 =B_{\mathrm{full}}-B_0^{\mathrm{fin}}.
\end{equation}
Маршрут остаётся живым только если
\begin{equation}
 B_{\mathrm{full}}>0
\end{equation}
в одной общей схеме вычитаний.

Машинный аудит сохранён в
\path{s2t_v3_finite_zero_mode_supertrace_results.json}.